% SEGMENT OLP-0721-S01
% Part: sets-functions-relations
% Chapter: sets
% Section: proofs-about-sets

\documentclass[../../../include/open-logic-section]{subfiles}

\begin{document}

\olfileid{sfr}{set}{prf}
\olsection{கணங்கள் பற்றிய நிறுவல்கள்}

\begin{editorial}
இப்பிரிவுக்குப் பதிலாக \verb|content/methods/proofs| பயன்படுத்தப்படுகிறது;
எனவே இயல்பான அமைப்பில் இப்பிரிவு இந்த அத்தியாயத்திலிருந்து
நீக்கப்பட்டுள்ளது.
\end{editorial}

% SEGMENT OLP-0721-S02
\begin{explain}
கணங்களும் இதுவரை அறிமுகப்படுத்திய குறியீடுகளும்,
கணங்களைக் குறிப்பிடவும் அவற்றுக்கிடையிலான தொடர்புகளை
வெளிப்படுத்தவும் வசதியான சுருக்கக் குறியீடுகளைத் தருகின்றன.
அத்தகைய தொடர்புகள் பற்றிய கூற்றுகளை நிறுவுவதும் அடிக்கடி
தேவைப்படும். கணித நிறுவல்களை நீங்கள் அறிந்திருக்காவிட்டால்,
இது புதிதாக இருக்கலாம். ஆகவே ஓர் எளிய எடுத்துக்காட்டைப்
படிப்படியாகப் பார்ப்போம். எந்தக் கணங்கள் $X$, $Y$ க்கும்
$X \cap (X \cup Y) = X$ எப்போதும் பொருந்தும் என்று
நிறுவுவோம். இத்தகைய கணச் சமத்துவத்தை எவ்வாறு நிறுவுவது?
கணங்கள் அவற்றின் !!{element}s ஆல் மட்டுமே தீர்மானிக்கப்படுகின்றன
என்பதை நினைவுகூர்க. அதாவது, கணங்கள் ஒரே !!{element}s ஐக்
கொண்டிருந்தால், அப்போதும் அப்போது மட்டுமே, அவை சமம்.
எனவே இங்கு (a) $X \cap (X \cup Y)$ இன் ஒவ்வோர் !!{element}
உம் $X$ இன் !!{element} என்றும், மறுதிசையில்
(b) $X$ இன் ஒவ்வோர் !!{element} உம் $X \cap (X \cup Y)$ இன்
!!{element} என்றும் நிறுவ வேண்டும். வேறு சொற்களில்,
(a) $X \cap (X \cup Y) \subseteq X$ என்பதையும்
(b) $X \subseteq X \cap (X \cup Y)$ என்பதையும் காட்டுகிறோம்.

% SEGMENT OLP-0721-S03
``$X \cap (X \cup Y)$ இன் ஒவ்வோர் !!{element}~$z$ உம்
$X$ இன் !!{element}'' போன்ற பொதுக் கூற்றை நிறுவும்போது,
முதலில் தன்னிச்சையான $z \in X \cap (X \cup Y)$ கொடுக்கப்பட்டுள்ளது
என்று எடுத்துக்கொண்டு, அந்த எடுகோளிலிருந்து $z \in X$ என்பதை
நிறுவுகிறோம். இவ்வடிவத்தை ``பொது நிபந்தனை நிறுவல்'' என்று
நீங்கள் அறிந்திருக்கலாம். இந்நிறுவலில், எடுகோளிலும் முடிவிலும்
பயன்படும் வரையறைகளையும் பயன்படுத்த வேண்டும்; இங்கு
``$\cap$'', ``$\cup$'' ஆகியவற்றின் வரையறைகள் அவை.
எனவே (a) நிகழ்வை பின்வருமாறு நிறுவலாம்:

% SEGMENT OLP-0721-S04
\begin{quote}
(a) முதலில் $X \cap (X \cup Y) \subseteq X$ என்பதைக்
காட்ட விரும்புகிறோம். அதாவது, $\subseteq$ இன் வரையறையின்படி,
எந்த~$z$ க்கும் $z \in X \cap (X \cup Y)$ எனில்
$z \in X$ என்று காட்ட வேண்டும். ஆகவே
$z \in X \cap (X \cup Y)$ என்க. $z$ என்ற !!{element},
இரண்டு கணங்களின் வெட்டுக் கணத்தில் இருப்பது, இரு கணங்களிலும்
!!{element} ஆக இருந்தால், அப்போதும் அப்போது மட்டுமே ஆகும்.
எனவே $z \in X$ என்றும் $z \in X \cup Y$ என்றும் முடிவு
செய்யலாம். குறிப்பாக, $z \in X$; இதையே காட்ட விரும்பினோம்.
\end{quote}

% SEGMENT OLP-0721-S05
இதனால் நிறுவலின் முதல் பாதி முடிந்தது. கடைசிப் படியில்,
ஓர் எடுகோளிலிருந்து ஓர் இணையல் ($z \in X$ மற்றும்
$z \in X \cup Y$) பின்தொடர்ந்தால், அதே எடுகோளிலிருந்து
அதன் ஒவ்வோர் இணைக் கூற்றும் பின்தொடரும் என்ற உண்மையைப்
பயன்படுத்தினோம். இவ்விதியை ``இணையல் நீக்கல்'' அல்லது
$\land$ நீக்கல் என்று நீங்கள் அறிந்திருக்கலாம்.
இப்போது (b) ஐ நிறுவுவோம்:

% SEGMENT OLP-0721-S06
\begin{quote}
(b) இப்போது $X \subseteq X \cap (X \cup Y)$ என்பதை
நிறுவுகிறோம். அதாவது, $\subseteq$ இன் வரையறையின்படி,
எந்த~$z$ க்கும் $z \in X$ எனில்
$z \in X \cap (X \cup Y)$ உம் ஆகும்.
$z \in X$ என்க. $z \in X \cap (X \cup Y)$ என்பதை
நிறுவ, ``$\cap$'' இன் வரையறையின்படி, (i) $z \in X$
என்பதையும் (ii) $z \in X \cup Y$ என்பதையும் காட்ட வேண்டும்.
(i) நமது எடுகோளே; ஆகவே மேலும் நிறுவ வேண்டியது ஏதுமில்லை.
(ii) க்கு, $z$, கணங்களின் சேர்ப்புக் கணத்தின் !!{element}
ஆக இருப்பது, அக்கணங்களில் குறைந்தது ஒன்றின் !!{element}
ஆக இருந்தால், அப்போதும் அப்போது மட்டுமே என்பதை நினைவுகூர்க.
$z \in X$ என்பதாலும், $X \cup Y$ என்பது $X$, $Y$ இன்
சேர்ப்புக் கணம் என்பதாலும், இங்கு அவ்வாறே உள்ளது.
ஆகவே $z \in X \cup Y$. (i) $z \in X$ என்பதையும்
(ii) $z \in X \cup Y$ என்பதையும் காட்டிவிட்டோம்;
எனவே ``$\cap$'' இன் வரையறையின்படி
$z \in X \cap (X \cup Y)$.
\end{quote}

% SEGMENT OLP-0721-S07
இது சற்றே விரிவான விளக்கம்; ஆனால் கணங்களையும்
அவற்றுக்கிடையிலான தொடர்புகளையும் பற்றி எவ்வாறு காரணம்
கூறுகிறோம் என்பதைக் காட்டுகிறது. வழக்கமாக இவ்வளவு
வெளிப்படையாக ஒவ்வொரு படியையும் கூறுவதில்லை; குறிப்பாக,
எல்லா வரையறைகளையும் மீண்டும் கூறாமல் இருக்கலாம்.
மேம்பட்ட பாடநூலில் இம்முடிவுக்கான நிறுவல் மிகவும்
சுருக்கமாக இருக்கும். அது பின்வருமாறு அமையலாம்.
\end{explain}

% SEGMENT OLP-0721-S08
\begin{prop}[உட்கவர் பண்பு]
எல்லாக் கணங்கள் $X$, $Y$ க்கும்,
\[
X \cap (X \cup Y) = X
\]
\end{prop}

\begin{proof}
(a) $z \in X \cap (X \cup Y)$ என்க. அப்போது $z \in X$;
ஆகவே $X \cap (X \cup Y) \subseteq X$.

(b) இப்போது $z \in X$ என்க. அப்போது $z \in X \cup Y$
உம், எனவே $z \in X \cap (X \cup Y)$ உம் ஆகும்.
ஆகவே $X \subseteq X \cap (X \cup Y)$.
\end{proof}

% SEGMENT OLP-0721-S09
\begin{prob}
$X \cup (X \cap Y) = X$ என்பதை விரிவாக நிறுவுக.
பின்னர் அதன் சுருக்கமான நிறுவலைத் தருக.
(குறிப்பு: $X \cup (X \cap Y) \subseteq X$ என்ற திசைக்கு,
நிகழ்வுகளைப் பிரித்து நிறுவுதல், அதாவது $\lor$ நீக்கல்,
தேவைப்படும்.)
\end{prob}

\end{document}
